% GENERATED FILE. Edit canonical lessons, then run npm run generate:books.
% canonical-source-hash: fnv1a64:f9bb186dd9013be5
% canonical-lessons: RU-C434-zhaloba, RU-C434-zhalovatsya, RU-C434-pretenziya, RU-C434-nedovolnyy, RU-C434-neispravnyy, RU-R434-review-427-434

\chapter{Making a complaint}
\label{ch:ru-making-a-complaint}
\hlchaptermodality{434}

\begin{chapteropening}
\textbf{By the end of this chapter:} \emph{I can make a complaint: the complaint, complaining, a formal claim, being dissatisfied and things that are out of order.}

Complete the last lesson of chapter 434: I can make a complaint: the complaint, complaining, a formal claim, being dissatisfied and things that are out of order.
\end{chapteropening}

\section[\texorpdfstring{\ru{жалоба} (zháloba)}{жалоба (zháloba)}]{\ru{жалоба} (zháloba) — a complaint}

\label{lesson:RU-C434-zhaloba}



\begin{quote}
Before the new one: say the Russian for a safe, then the Russian for a fault.
\end{quote}

\subsection*{You'll want to know: \ru{жалоба}}
\textbf{\ru{жалоба}} (\emph{zháloba}) — \textquotedblleft{}a complaint\textquotedblright{}.

\begin{quote}
\ru{Туристы} \ru{подали} \ru{жалобу} \ru{на} \ru{шум}.
\end{quote}

\begin{quote}
\emph{Turísty podáli zhálobu na shum.}
\end{quote}

\begin{quote}
The tourists made a complaint about the noise.
\end{quote}

Feminine. Takes \emph{\ru{на}} + accusative: \emph{\ru{жалоба} \ru{на} \ru{соседей}}.

\begin{cousinweb}
\emph{\ru{жалоба}} goes with \emph{\ru{жалеть}} 'to pity' and \emph{\ru{жаль}} 'a pity': to complain is to call for pity. EtymDB lists \emph{\ru{жалоба}} with the sense 'complaint' but without a source.
\end{cousinweb}

\begin{grammarlens}[title={What you've built}]
One more word for telling, explaining and describing at B1.
\end{grammarlens}

\subsection*{Your turn}
\begin{itemize}
\raggedright
  \item \emph{Say it:} \emph{zháloba}
  \item \emph{Say it:} \emph{zháloba}, once more
  \item \emph{Say it:} a sentence of your own with \emph{zháloba}
  \item \emph{Recall:} say the Russian for a safe, then the Russian for a fault, then say \emph{zháloba} again
\end{itemize}

\begin{tcolorbox}[breakable,colback=teal!4,colframe=teal!35!black,title={Before you move on}]
What does \ru{жалоба} mean? (\textquotedblleft{}A complaint\textquotedblright{}.) Say it once more.
\end{tcolorbox}
\section[\texorpdfstring{\ru{жаловаться} (zhálovatsya)}{жаловаться (zhálovatsya)}]{\ru{жаловаться} (zhálovatsya) — to complain}

\label{lesson:RU-C434-zhalovatsya}



\begin{quote}
Before the new one: say the Russian for a fault, then the Russian for a complaint.
\end{quote}

\subsection*{You'll want to know: \ru{жаловаться}}
\textbf{\ru{жаловаться}} (\emph{zhálovatsya}) — \textquotedblleft{}to complain\textquotedblright{}.

\begin{quote}
\ru{Дед} \ru{всё} \ru{время} \ru{жалуется} \ru{на} \ru{погоду}.
\end{quote}

\begin{quote}
\emph{Ded vsyo vrémya zháluyetsya na pogódu.}
\end{quote}

\begin{quote}
Grandad complains about the weather all the time.
\end{quote}

Imperfective; the perfective is \emph{\ru{пожаловаться}}. Present \emph{\ru{жалуюсь}, \ru{жалуешься}}. \emph{\ru{На}} + accusative for what, the dative for whom.

\begin{cousinweb}
\emph{\ru{жаловаться}} goes with \emph{\ru{жалоба}} 'a complaint' and \emph{\ru{жалеть}} 'to pity'. EtymDB lists \emph{\ru{жалоба}} and \emph{\ru{жалеть}} without a deeper source.
\end{cousinweb}

\begin{grammarlens}[title={What you've built}]
One more word for telling, explaining and describing at B1.
\end{grammarlens}

\subsection*{Your turn}
\begin{itemize}
\raggedright
  \item \emph{Say it:} \emph{zhálovatsya}
  \item \emph{Say it:} \emph{zhálovatsya}, once more
  \item \emph{Say it:} a sentence of your own with \emph{zhálovatsya}
  \item \emph{Recall:} say the Russian for a fault, then the Russian for a complaint, then say \emph{zhálovatsya} again
\end{itemize}

\begin{tcolorbox}[breakable,colback=teal!4,colframe=teal!35!black,title={Before you move on}]
What does \ru{жаловаться} mean? (\textquotedblleft{}To complain\textquotedblright{}.) Say it once more.
\end{tcolorbox}
\section[\texorpdfstring{\ru{претензия} (preténziya)}{претензия (preténziya)}]{\ru{претензия} (preténziya) — a formal complaint, a claim}

\label{lesson:RU-C434-pretenziya}



\begin{quote}
Before the new one: say the Russian for a complaint, then the Russian for to complain.
\end{quote}

\subsection*{You'll want to know: \ru{претензия}}
\textbf{\ru{претензия}} (\emph{preténziya}) — \textquotedblleft{}a formal complaint, a claim\textquotedblright{}.

\begin{quote}
\ru{Мы} \ru{отправили} \ru{претензию} \ru{в} \ru{турфирму}.
\end{quote}

\begin{quote}
\emph{My otprávili preténziyu v turfírmu.}
\end{quote}

\begin{quote}
We sent a formal complaint to the travel agency.
\end{quote}

Feminine. \emph{\ru{У} \ru{меня} \ru{к} \ru{вам} \ru{претензия}} — I have a complaint to make to you.

\begin{cousinweb}
\emph{\ru{претензия}} is usually traced, probably through German, to Medieval Latin \emph{praetēnsiō} 'a claim', from \emph{praetendere} 'to stretch out, to put forward' (Lewis \& Short). EtymDB does not list it, so the route is hedged.
\end{cousinweb}

\begin{grammarlens}[title={What you've built}]
One more word for telling, explaining and describing at B1.
\end{grammarlens}

\subsection*{Your turn}
\begin{itemize}
\raggedright
  \item \emph{Say it:} \emph{preténziya}
  \item \emph{Say it:} \emph{preténziya}, once more
  \item \emph{Say it:} a sentence of your own with \emph{preténziya}
  \item \emph{Recall:} say the Russian for a complaint, then the Russian for to complain, then say \emph{preténziya} again
\end{itemize}

\begin{tcolorbox}[breakable,colback=teal!4,colframe=teal!35!black,title={Before you move on}]
What does \ru{претензия} mean? (\textquotedblleft{}A formal complaint, a claim\textquotedblright{}.) Say it once more.
\end{tcolorbox}
\section[\texorpdfstring{\ru{недовольный} (nedovólnyy)}{недовольный (nedovólnyy)}]{\ru{недовольный} (nedovólnyy) — dissatisfied, displeased}

\label{lesson:RU-C434-nedovolnyy}



\begin{quote}
Before the new one: say the Russian for to complain, then the Russian for a formal complaint.
\end{quote}

\subsection*{You'll want to know: \ru{недовольный}}
\textbf{\ru{недовольный}} (\emph{nedovólnyy}) — \textquotedblleft{}dissatisfied, displeased\textquotedblright{}.

\begin{quote}
\ru{Недовольный} \ru{клиент} \ru{долго} \ru{спорил} \ru{с} \ru{кассиром}.
\end{quote}

\begin{quote}
\emph{Nedovólnyy kliyént dólgo spóril s kassírom.}
\end{quote}

\begin{quote}
The dissatisfied customer argued with the cashier for a long time.
\end{quote}

Takes the instrumental: \emph{\ru{недоволен} \ru{номером}} — unhappy with the room.

\begin{cousinweb}
\emph{\ru{недовольный}} is built from \emph{\ru{не}-} 'not' and \emph{\ru{довольный}} 'content, satisfied' (EtymDB lists both); \emph{\ru{довольный}} goes with \emph{\ru{довольно}} 'enough'.
\end{cousinweb}

\begin{grammarlens}[title={What you've built}]
One more word for telling, explaining and describing at B1.
\end{grammarlens}

\subsection*{Your turn}
\begin{itemize}
\raggedright
  \item \emph{Say it:} \emph{nedovólnyy}
  \item \emph{Say it:} \emph{nedovólnyy}, once more
  \item \emph{Say it:} a sentence of your own with \emph{nedovólnyy}
  \item \emph{Recall:} say the Russian for to complain, then the Russian for a formal complaint, then say \emph{nedovólnyy} again
\end{itemize}

\begin{tcolorbox}[breakable,colback=teal!4,colframe=teal!35!black,title={Before you move on}]
What does \ru{недовольный} mean? (\textquotedblleft{}Dissatisfied, displeased\textquotedblright{}.) Say it once more.
\end{tcolorbox}
\section[\texorpdfstring{\ru{неисправный} (neisprávnyy)}{неисправный (neisprávnyy)}]{\ru{неисправный} (neisprávnyy) — out of order, faulty}

\label{lesson:RU-C434-neispravnyy}



\begin{quote}
Before the new one: say the Russian for a formal complaint, then the Russian for dissatisfied.
\end{quote}

\subsection*{You'll want to know: \ru{неисправный}}
\textbf{\ru{неисправный}} (\emph{neisprávnyy}) — \textquotedblleft{}out of order, faulty\textquotedblright{}.

\begin{quote}
\ru{В} \ru{номере} \ru{был} \ru{неисправный} \ru{холодильник}.
\end{quote}

\begin{quote}
\emph{V nómere byl neisprávnyy kholodílnik.}
\end{quote}

\begin{quote}
There was a fridge in the room that was out of order.
\end{quote}

The opposite is \emph{\ru{исправный}} — in working order.

\begin{cousinweb}
\emph{\ru{неисправный}} is built from \emph{\ru{не}-} 'not' and \emph{\ru{исправный}} 'in good order', from \emph{\ru{исправить}} 'to put right', on \emph{\ru{править}} 'to set right'. EtymDB lists \emph{\ru{неисправный}} without a deeper source.
\end{cousinweb}

\begin{grammarlens}[title={What you've built}]
One more word for telling, explaining and describing at B1.
\end{grammarlens}

\subsection*{Your turn}
\begin{itemize}
\raggedright
  \item \emph{Say it:} \emph{neisprávnyy}
  \item \emph{Say it:} \emph{neisprávnyy}, once more
  \item \emph{Say it:} a sentence of your own with \emph{neisprávnyy}
  \item \emph{Recall:} say the Russian for a formal complaint, then the Russian for dissatisfied, then say \emph{neisprávnyy} again
\end{itemize}

\begin{tcolorbox}[breakable,colback=teal!4,colframe=teal!35!black,title={Before you move on}]
What does \ru{неисправный} mean? (\textquotedblleft{}Out of order, faulty\textquotedblright{}.) Say it once more.
\end{tcolorbox}
\section[(dialogue)]{Words from chapters 427–434 — a review}

\label{lesson:RU-R434-review-427-434}



\begin{quote}
Say the Russian for a harbour and a lighthouse.
\end{quote}

\subsection*{Your turn}
Say these aloud:

\begin{itemize}
\raggedright
  \item a harbour, a lighthouse, a sailor, a captain, to sink
  \item a hostel, a holiday camp, a bed for the night, to spend the night, a booking
  \item a receptionist, to check into a hotel, to check out, a deposit, to extend
  \item bedding, a sheet, a pillowcase, a chambermaid, to make up
  \item an air conditioner, a heater, a draught, to air, to dry
  \item to leak, to drip, to get blocked, to flow, a flood
  \item noise, to disturb, a cockroach, a safe, a fault
  \item a complaint, to complain, a formal complaint, dissatisfied, out of order
\end{itemize}

\begin{tcolorbox}[breakable,colback=teal!4,colframe=teal!35!black,title={Before you move on}]
A harbour? (\textbf{\ru{гавань}}.) A lighthouse? (\textbf{\ru{маяк}}.) A sailor? (\textbf{\ru{матрос}}.) A captain? (\textbf{\ru{капитан}}.) To sink? (\textbf{\ru{тонуть}}.) A hostel? (\textbf{\ru{хостел}}.) A holiday camp? (\textbf{\ru{турбаза}}.) A bed for the night? (\textbf{\ru{ночлег}}.) To spend the night? (\textbf{\ru{ночевать}}.) A booking? (\textbf{\ru{бронь}}.) A receptionist? (\textbf{\ru{администратор}}.) To check into a hotel? (\textbf{\ru{заселяться}}.) To check out? (\textbf{\ru{выселяться}}.) A deposit? (\textbf{\ru{залог}}.) To extend? (\textbf{\ru{продлевать}}.) Bedding? (\textbf{\ru{постель}}.) A sheet? (\textbf{\ru{простыня}}.) A pillowcase? (\textbf{\ru{наволочка}}.) A chambermaid? (\textbf{\ru{горничная}}.) To make up? (\textbf{\ru{стелить}}.) An air conditioner? (\textbf{\ru{кондиционер}}.) A heater? (\textbf{\ru{обогреватель}}.) A draught? (\textbf{\ru{сквозняк}}.) To air? (\textbf{\ru{проветривать}}.) To dry? (\textbf{\ru{сушить}}.) To leak? (\textbf{\ru{протекать}}.) To drip? (\textbf{\ru{капать}}.) To get blocked? (\textbf{\ru{засоряться}}.) To flow? (\textbf{\ru{течь}}.) A flood? (\textbf{\ru{потоп}}.) Noise? (\textbf{\ru{шум}}.) To disturb? (\textbf{\ru{мешать}}.) A cockroach? (\textbf{\ru{таракан}}.) A safe? (\textbf{\ru{сейф}}.) A fault? (\textbf{\ru{неполадка}}.) A complaint? (\textbf{\ru{жалоба}}.) To complain? (\textbf{\ru{жаловаться}}.) A formal complaint? (\textbf{\ru{претензия}}.) Dissatisfied? (\textbf{\ru{недовольный}}.) Out of order? (\textbf{\ru{неисправный}}.)
\end{tcolorbox}
